% Apercu autonome des tables exportees (etape 8). Genere par src/export, ne pas editer.
% Usage : importer le depot a la racine dans Overleaf, choisir ce fichier comme document
% principal, compiler UNE fois et consulter le journal (aucune compilation locale).
% Si le projet Overleaf n'a pas le depot a sa racine, modifier uniquement \tablesroot.
\documentclass[a4paper,10pt]{article}
\usepackage[T1]{fontenc}
\usepackage[utf8]{inputenc}
\usepackage{lmodern}
\usepackage[a4paper,margin=2cm]{geometry}
\usepackage{amsmath}
\usepackage{booktabs}
\usepackage{longtable}
\usepackage{adjustbox}
\usepackage{caption}
\usepackage{siunitx}
\sisetup{output-decimal-marker={,}}% presentation francaise ; les donnees gardent le point
\newcommand{\tablesroot}{results/tables}
\begin{document}
\section*{Aper\c{c}u des tables export\'ees}

\subsection*{\texttt{nominal\_1d\_seed42}}
Profil d'erreur nominal (graine 42), degr\'es 1 \`a 12.

\captionof{table}{nominal\_1d\_seed42}
\begingroup\footnotesize
\begin{adjustbox}{max width=\linewidth}
\input{\tablesroot/nominal_1d_seed42.tex}
\end{adjustbox}
\endgroup
\medskip

\subsection*{\texttt{nominal\_1d\_seed42\_diagnostics}}
Diagnostics de conditionnement et \'ecarts OLS/QR/SVD (graine 42).

\captionof{table}{nominal\_1d\_seed42\_diagnostics}
\begingroup\footnotesize
\begin{adjustbox}{max width=\linewidth}
\input{\tablesroot/nominal_1d_seed42_diagnostics.tex}
\end{adjustbox}
\endgroup
\medskip

\subsection*{\texttt{nominal\_1d\_stats\_50\_graines}}
Statistiques descriptives sur 50 graines (m\'ediane, quartiles).

\captionof{table}{nominal\_1d\_stats\_50\_graines}
\begingroup\footnotesize
\begin{adjustbox}{max width=\linewidth}
\input{\tablesroot/nominal_1d_stats_50_graines.tex}
\end{adjustbox}
\endgroup
\medskip

\subsection*{\texttt{nominal\_1d\_cstar\_50\_graines}}
Distribution de C* sur 50 graines.

\captionof{table}{nominal\_1d\_cstar\_50\_graines}
\begingroup\footnotesize
\begin{adjustbox}{max width=\linewidth}
\input{\tablesroot/nominal_1d_cstar_50_graines.tex}
\end{adjustbox}
\endgroup
\medskip

\subsection*{\texttt{exploration\_nb\_ex\_synthese}}
Exploration 1 (taille de l'\'echantillon N) : synth\`ese par N.

\captionof{table}{exploration\_nb\_ex\_synthese}
\begingroup\footnotesize
\begin{adjustbox}{max width=\linewidth}
\input{\tablesroot/exploration_nb_ex_synthese.tex}
\end{adjustbox}
\endgroup
\medskip

\subsection*{\texttt{exploration\_sigma\_synthese}}
Exploration 2 (bruit sigma) : synth\`ese par sigma.

\captionof{table}{exploration\_sigma\_synthese}
\begingroup\footnotesize
\begin{adjustbox}{max width=\linewidth}
\input{\tablesroot/exploration_sigma_synthese.tex}
\end{adjustbox}
\endgroup
\medskip

\subsection*{\texttt{exploration\_extrapolation\_stats}}
Exploration 3 (extrapolation) : EQM par domaine.

\captionof{table}{exploration\_extrapolation\_stats}
\begingroup\footnotesize
\begin{adjustbox}{max width=\linewidth}
\input{\tablesroot/exploration_extrapolation_stats.tex}
\end{adjustbox}
\endgroup
\medskip

\subsection*{\texttt{exploration\_nb\_ex\_stats}}
Exploration 1 : statistiques compl\`etes par N et par degr\'e.

\captionof{table}{exploration\_nb\_ex\_stats}
\begingroup\footnotesize
\input{\tablesroot/exploration_nb_ex_stats.tex}
\endgroup
\medskip

\subsection*{\texttt{exploration\_sigma\_stats}}
Exploration 2 : statistiques compl\`etes par sigma et par degr\'e.

\captionof{table}{exploration\_sigma\_stats}
\begingroup\footnotesize
\input{\tablesroot/exploration_sigma_stats.tex}
\endgroup
\medskip

\subsection*{\texttt{exploration\_nb\_ex\_diagnostics}}
Diagnostics de conditionnement par N et par degr\'e.

\captionof{table}{exploration\_nb\_ex\_diagnostics}
\begingroup\footnotesize
\input{\tablesroot/exploration_nb_ex_diagnostics.tex}
\endgroup
\medskip

\subsection*{\texttt{exploration\_nb\_ex\_variance}}
Variance de pr\'ediction moyenne par N (\'echelle en 1/N).

\captionof{table}{exploration\_nb\_ex\_variance}
\begingroup\footnotesize
\input{\tablesroot/exploration_nb_ex_variance.tex}
\endgroup
\medskip

\subsection*{\texttt{mlp\_2d\_nominal\_stats}}
EXP-2D-NOM : EQM de validation par largeur H, 10 graines.

\captionof{table}{mlp\_2d\_nominal\_stats}
\begingroup\footnotesize
\begin{adjustbox}{max width=\linewidth}
\input{\tablesroot/mlp_2d_nominal_stats.tex}
\end{adjustbox}
\endgroup
\medskip

\subsection*{\texttt{mlp\_2d\_sensibilite\_topologie}}
EXP-2D-TOPO : profondeur contre largeur.

\captionof{table}{mlp\_2d\_sensibilite\_topologie}
\begingroup\footnotesize
\begin{adjustbox}{max width=\linewidth}
\input{\tablesroot/mlp_2d_sensibilite_topologie.tex}
\end{adjustbox}
\endgroup
\medskip

\subsection*{\texttt{mlp\_2d\_sensibilite\_regularisation}}
EXP-2D-REG : r\'egularisation L2.

\captionof{table}{mlp\_2d\_sensibilite\_regularisation}
\begingroup\footnotesize
\begin{adjustbox}{max width=\linewidth}
\input{\tablesroot/mlp_2d_sensibilite_regularisation.tex}
\end{adjustbox}
\endgroup
\medskip

\subsection*{\texttt{mlp\_2d\_sensibilite\_activation}}
EXP-2D-ACT : tanh contre relu.

\captionof{table}{mlp\_2d\_sensibilite\_activation}
\begingroup\footnotesize
\begin{adjustbox}{max width=\linewidth}
\input{\tablesroot/mlp_2d_sensibilite_activation.tex}
\end{adjustbox}
\endgroup
\medskip

\subsection*{\texttt{mlp\_2d\_sensibilite\_solveur}}
EXP-2D-SOLV : lbfgs contre adam (budget et grille donn\'es).

\captionof{table}{mlp\_2d\_sensibilite\_solveur}
\begingroup\footnotesize
\begin{adjustbox}{max width=\linewidth}
\input{\tablesroot/mlp_2d_sensibilite_solveur.tex}
\end{adjustbox}
\endgroup
\medskip

\subsection*{\texttt{mlp\_2d\_adam\_grid}}
Grille adam (DEC-19), ordre contractuel de la grille.

\captionof{table}{mlp\_2d\_adam\_grid}
\begingroup\footnotesize
\begin{adjustbox}{max width=\linewidth}
\input{\tablesroot/mlp_2d_adam_grid.tex}
\end{adjustbox}
\endgroup
\medskip

\subsection*{\texttt{mlp\_2d\_solver\_comparison}}
Comparaison lbfgs / adam sur le test (10 graines).

\captionof{table}{mlp\_2d\_solver\_comparison}
\begingroup\footnotesize
\begin{adjustbox}{max width=\linewidth}
\input{\tablesroot/mlp_2d_solver_comparison.tex}
\end{adjustbox}
\endgroup
\medskip

\subsection*{\texttt{mlp\_2d\_final\_synthese}}
\'Evaluation finale sur le test (N = 1600), 10 graines par configuration.

\captionof{table}{mlp\_2d\_final\_synthese}
\begingroup\footnotesize
\begin{adjustbox}{max width=\linewidth}
\input{\tablesroot/mlp_2d_final_synthese.tex}
\end{adjustbox}
\endgroup
\medskip

\subsection*{\texttt{mlp\_2d\_final\_test\_stats}}
\'Evaluation finale : toutes les m\'etriques.

\captionof{table}{mlp\_2d\_final\_test\_stats}
\begingroup\footnotesize
\begin{adjustbox}{max width=\linewidth}
\input{\tablesroot/mlp_2d_final_test_stats.tex}
\end{adjustbox}
\endgroup
\medskip

\subsection*{\texttt{mlp\_2d\_final\_test\_runs}}
\'Evaluation finale : r\'esultat de chaque graine.

\captionof{table}{mlp\_2d\_final\_test\_runs}
\begingroup\footnotesize
\input{\tablesroot/mlp_2d_final_test_runs.tex}
\endgroup
\medskip

\subsection*{\texttt{mlp\_2d\_configurations}}
Toutes les configurations distinctes du plan (validation).

\captionof{table}{mlp\_2d\_configurations}
\begingroup\footnotesize
\begin{adjustbox}{max width=\linewidth}
\input{\tablesroot/mlp_2d_configurations.tex}
\end{adjustbox}
\endgroup
\medskip

\subsection*{\texttt{donnees\_2d\_provenance}}
Provenance et empreintes des donn\'ees 2D.

\captionof{table}{donnees\_2d\_provenance}
\begingroup\footnotesize
\begin{adjustbox}{max width=\linewidth}
\input{\tablesroot/donnees_2d_provenance.tex}
\end{adjustbox}
\endgroup
\medskip

\end{document}
